% GENERATED FILE. Edit canonical lessons, then run npm run generate:books.
% canonical-source-hash: fnv1a64:cc76017b6c879fe8
% canonical-lessons: ML-C565-citrakaran, ML-C565-samgitajnan, ML-C565-narttaki, ML-C565-eluttukaran, ML-C565-nirmmikkuka

\chapter{Painters, Dancers and Writers}
\label{ch:ml-painters-dancers-and-writers}
\hlchaptermodality{565}

\begin{chapteropening}
\textbf{By the end of this chapter:} \emph{I can say a painter, a musician, a dancer, a writer and to produce a film.}

Complete the last lesson of chapter 565: I can say a painter, a musician, a dancer, a writer and to produce a film.
\end{chapteropening}

\section[\texorpdfstring{\ml{ചിത്രകാരൻ} (citrakāran)}{ചിത്രകാരൻ (citrakāran)}]{\ml{ചിത്രകാരൻ} (citrakāran) — a painter}

\label{lesson:ML-C565-citrakaran}



\begin{quote}
Before the new one: say the Malayalam for to strum, then the Malayalam for to train in.
\end{quote}

\subsection*{You'll want to know: \ml{ചിത്രകാരൻ}}
\textbf{\ml{ചിത്രകാരൻ}} — \emph{citrakāran} — \textquotedblleft{}a painter\textquotedblright{}.

\begin{quote}
\ml{ആ} \ml{ചിത്രകാരന്റെ} \ml{ചിത്രങ്ങൾ} \ml{പ്രദർശനത്തിൽ} \ml{ഉണ്ടായിരുന്നു}.
\end{quote}

\begin{quote}
\emph{ā citrakāranṟe citraṅṅaḷ pradarśanattil uṇṭāyirunnu.}
\end{quote}

\begin{quote}
That painter's pictures were in the exhibition.
\end{quote}

A woman painter is \ml{ചിത്രകാരി}.

\begin{cousinweb}
\textbf{\ml{ചിത്രകാരൻ}} (\emph{citrakāran}) is borrowed from Sanskrit \emph{citrakāra}, \textquotedblleft{}a painter\textquotedblright{} (CDIAL 4807), \emph{citra}, \textquotedblleft{}a picture\textquotedblright{}, and \emph{kāra}, \textquotedblleft{}a maker\textquotedblright{}.
\end{cousinweb}

\begin{grammarlens}[title={What you've built}]
One more word for describing what you lived and felt.
\end{grammarlens}

\subsection*{Your turn}
\begin{itemize}
\raggedright
  \item \emph{Say it:} \emph{citrakāran}
  \item \emph{Say it:} \emph{citrakāran}, once more
  \item \emph{Say it:} \emph{ā citrakāranṟe citraṅṅaḷ pradarśanattil uṇṭāyirunnu.} — That painter's pictures were in the exhibition.
  \item \emph{Recall:} say the Malayalam for to strum, then the Malayalam for to train in, then say \emph{citrakāran} again
\end{itemize}

\begin{tcolorbox}[breakable,colback=teal!4,colframe=teal!35!black,title={Before you move on}]
What does \ml{ചിത്രകാരൻ} mean? (\textquotedblleft{}A painter\textquotedblright{}.) Say it once more.
\end{tcolorbox}
\section[\texorpdfstring{\ml{സംഗീതജ്ഞൻ} (saṁgītajñan)}{സംഗീതജ്ഞൻ (saṁgītajñan)}]{\ml{സംഗീതജ്ഞൻ} (saṁgītajñan) — a musician}

\label{lesson:ML-C565-samgitajnan}



\begin{quote}
Before the new one: say the Malayalam for to train in, then the Malayalam for a painter.
\end{quote}

\subsection*{You'll want to know: \ml{സംഗീതജ്ഞൻ}}
\textbf{\ml{സംഗീതജ്ഞൻ}} — \emph{saṁgītajñan} — \textquotedblleft{}a musician\textquotedblright{}.

\begin{quote}
\ml{ഒരു} \ml{പ്രശസ്ത} \ml{സംഗീതജ്ഞൻ} \ml{കച്ചേരിയിൽ} \ml{പാടി}.
\end{quote}

\begin{quote}
\emph{oru praśasta saṁgītajñan kaccēriyil pāṭi.}
\end{quote}

\begin{quote}
A famous musician sang at the concert.
\end{quote}

\ml{സംഗീതം} music + -\ml{ജ്ഞൻ}, one who knows, as in \ml{ശാസ്ത്രജ്ഞൻ}.

\begin{cousinweb}
\textbf{\ml{സംഗീതജ്ഞൻ}} (\emph{saṁgītajñan}) is a learned compound built from \ml{സംഗീതം}, from Sanskrit \emph{saṁgīta}, \textquotedblleft{}a chorus, a concert\textquotedblright{} (CDIAL 12849), and \emph{-jña}, \textquotedblleft{}knowing\textquotedblright{}.
\end{cousinweb}

\begin{grammarlens}[title={What you've built}]
One more word for describing what you lived and felt.
\end{grammarlens}

\subsection*{Your turn}
\begin{itemize}
\raggedright
  \item \emph{Say it:} \emph{saṁgītajñan}
  \item \emph{Say it:} \emph{saṁgītajñan}, once more
  \item \emph{Say it:} \emph{oru praśasta saṁgītajñan kaccēriyil pāṭi.} — A famous musician sang at the concert.
  \item \emph{Recall:} say the Malayalam for to train in, then the Malayalam for a painter, then say \emph{saṁgītajñan} again
\end{itemize}

\begin{tcolorbox}[breakable,colback=teal!4,colframe=teal!35!black,title={Before you move on}]
What does \ml{സംഗീതജ്ഞൻ} mean? (\textquotedblleft{}A musician\textquotedblright{}.) Say it once more.
\end{tcolorbox}
\section[\texorpdfstring{\ml{നർത്തകി} (narttaki)}{നർത്തകി (narttaki)}]{\ml{നർത്തകി} (narttaki) — a dancer}

\label{lesson:ML-C565-narttaki}



\begin{quote}
Before the new one: say the Malayalam for a painter, then the Malayalam for a musician.
\end{quote}

\subsection*{You'll want to know: \ml{നർത്തകി}}
\textbf{\ml{നർത്തകി}} — \emph{narttaki} — \textquotedblleft{}a dancer\textquotedblright{}.

\begin{quote}
\ml{ആ} \ml{നർത്തകി} \ml{വേദിയിൽ} \ml{മനോഹരമായി} \ml{നൃത്തം} \ml{ചെയ്തു}.
\end{quote}

\begin{quote}
\emph{ā narttaki vēdiyil manōharamāyi nṛttaṁ ceytu.}
\end{quote}

\begin{quote}
That dancer danced beautifully on the stage.
\end{quote}

A male dancer is \ml{നർത്തകൻ}.

\begin{cousinweb}
\textbf{\ml{നർത്തകി}} (\emph{narttaki}) is borrowed from Sanskrit \emph{nartakī}, \textquotedblleft{}a female dancer, an actress\textquotedblright{} (Monier-Williams), from the root \emph{nṛt}, \textquotedblleft{}to dance\textquotedblright{}, the root of \ml{നൃത്തം}.
\end{cousinweb}

\begin{grammarlens}[title={What you've built}]
One more word for describing what you lived and felt.
\end{grammarlens}

\subsection*{Your turn}
\begin{itemize}
\raggedright
  \item \emph{Say it:} \emph{narttaki}
  \item \emph{Say it:} \emph{narttaki}, once more
  \item \emph{Say it:} \emph{ā narttaki vēdiyil manōharamāyi nṛttaṁ ceytu.} — That dancer danced beautifully on the stage.
  \item \emph{Recall:} say the Malayalam for a painter, then the Malayalam for a musician, then say \emph{narttaki} again
\end{itemize}

\begin{tcolorbox}[breakable,colback=teal!4,colframe=teal!35!black,title={Before you move on}]
What does \ml{നർത്തകി} mean? (\textquotedblleft{}A dancer\textquotedblright{}.) Say it once more.
\end{tcolorbox}
\section[\texorpdfstring{\ml{എഴുത്തുകാരൻ} (eḻuttukāran)}{എഴുത്തുകാരൻ (eḻuttukāran)}]{\ml{എഴുത്തുകാരൻ} (eḻuttukāran) — a writer; an author}

\label{lesson:ML-C565-eluttukaran}



\begin{quote}
Before the new one: say the Malayalam for a musician, then the Malayalam for a dancer.
\end{quote}

\subsection*{You'll want to know: \ml{എഴുത്തുകാരൻ}}
\textbf{\ml{എഴുത്തുകാരൻ}} — \emph{eḻuttukāran} — \textquotedblleft{}a writer; an author\textquotedblright{}.

\begin{quote}
\ml{എന്റെ} \ml{ഇഷ്ട} \ml{എഴുത്തുകാരന്റെ} \ml{പുതിയ} \ml{നോവൽ} \ml{ഇറങ്ങി}.
\end{quote}

\begin{quote}
\emph{enṟe iṣṭa eḻuttukāranṟe putiya nōval iṟaṅṅi.}
\end{quote}

\begin{quote}
My favourite writer's new novel has come out.
\end{quote}

\ml{എഴുത്ത്} writing + -\ml{കാരൻ}, one who does it; a woman writer is \ml{എഴുത്തുകാരി}.

\begin{cousinweb}
\textbf{\ml{എഴുത്തുകാരൻ}} (\emph{eḻuttukāran}) is native Dravidian, built from \ml{എഴുത്ത്}, \textquotedblleft{}writing, a letter\textquotedblright{} (DEDR 853), with the ending -\ml{കാരൻ}, \textquotedblleft{}one who does it\textquotedblright{}.
\end{cousinweb}

\begin{grammarlens}[title={What you've built}]
One more word for describing what you lived and felt.
\end{grammarlens}

\subsection*{Your turn}
\begin{itemize}
\raggedright
  \item \emph{Say it:} \emph{eḻuttukāran}
  \item \emph{Say it:} \emph{eḻuttukāran}, once more
  \item \emph{Say it:} \emph{enṟe iṣṭa eḻuttukāranṟe putiya nōval iṟaṅṅi.} — My favourite writer's new novel has come out.
  \item \emph{Recall:} say the Malayalam for a musician, then the Malayalam for a dancer, then say \emph{eḻuttukāran} again
\end{itemize}

\begin{tcolorbox}[breakable,colback=teal!4,colframe=teal!35!black,title={Before you move on}]
What does \ml{എഴുത്തുകാരൻ} mean? (\textquotedblleft{}A writer; an author\textquotedblright{}.) Say it once more.
\end{tcolorbox}
\section[\texorpdfstring{\ml{നിർമ്മിക്കുക} (nirmmikkuka)}{നിർമ്മിക്കുക (nirmmikkuka)}]{\ml{നിർമ്മിക്കുക} (nirmmikkuka) — to produce a film; to make}

\label{lesson:ML-C565-nirmmikkuka}



\begin{quote}
Before the new one: say the Malayalam for a dancer, then the Malayalam for a writer.
\end{quote}

\subsection*{You'll want to know: \ml{നിർമ്മിക്കുക}}
\textbf{\ml{നിർമ്മിക്കുക}} — \emph{nirmmikkuka} — \textquotedblleft{}to produce a film; to make\textquotedblright{}.

\begin{quote}
\ml{ആ} \ml{കമ്പനി} \ml{ഒരുപാട്} \ml{സിനിമകൾ} \ml{നിർമ്മിച്ചു}.
\end{quote}

\begin{quote}
\emph{ā kampani orupāṭŭ sinimakaḷ nirmmiccu.}
\end{quote}

\begin{quote}
That company produced many films.
\end{quote}

Past: \ml{നിർമ്മിച്ചു}, \emph{nirmmiccu}.

To produce a film, to build a road or a bridge; a producer is \ml{നിർമ്മാതാവ്}.

\begin{cousinweb}
\textbf{\ml{നിർമ്മിക്കുക}} (\emph{nirmmikkuka}) is built from Sanskrit \emph{nirmāṇa}, \textquotedblleft{}measuring, making\textquotedblright{} (CDIAL 7368), with the Malayalam verb ending -\ml{ിക്കുക}.
\end{cousinweb}

\begin{grammarlens}[title={What you've built}]
One more word for describing what you lived and felt.
\end{grammarlens}

\subsection*{Your turn}
\begin{itemize}
\raggedright
  \item \emph{Say it:} \emph{nirmmikkuka}
  \item \emph{Say it:} \emph{nirmmikkuka}, once more
  \item \emph{Say it:} \emph{ā kampani orupāṭŭ sinimakaḷ nirmmiccu.} — That company produced many films.
  \item \emph{Recall:} say the Malayalam for a dancer, then the Malayalam for a writer, then say \emph{nirmmikkuka} again
\end{itemize}

\begin{tcolorbox}[breakable,colback=teal!4,colframe=teal!35!black,title={Before you move on}]
What does \ml{നിർമ്മിക്കുക} mean? (\textquotedblleft{}To produce a film; to make\textquotedblright{}.) Say it once more.
\end{tcolorbox}
